\documentclass[a4paper,UKenglish,cleveref,autoref,anonymous]{oasics-v2021}
% Anonymous working outline; not a submission-ready paper.
% Keep the publisher class unchanged. Planning and evidence status: README.md.
\hideOASIcs
\bibliographystyle{plainurl}

\title{CrossStage-Stitch: A Scene-Aware Runtime for Quality-Constrained GPU Panorama Stitching}
\titlerunning{CrossStage-Stitch: Quality-Constrained GPU Panorama Stitching}
\author{Anonymous Authors}{Anonymous affiliation}{}{}{}
\authorrunning{Anonymous Authors}
\Copyright{Anonymous Authors}
\ccsdesc[500]{Computing methodologies~Parallel computing methodologies}
\keywords{GPU panorama stitching, adaptive computation, runtime management, quality of result, cross-stage reuse}
\category{}
% Publication identifiers and proceedings editors are assigned by the publisher.
% Do not copy the sample's fictitious event, volume, article number, or DOI.
\AtBeginDocument{\hypersetup{pdfsubject={Anonymous working outline for PARMA-DITAM 2027}}}

\begin{document}
\maketitle

\begin{abstract}
% Role: scope and research question; empirical findings are not yet available.
This working outline defines a study of scene-aware computational budgets in a
GPU panorama-stitching pipeline. The research question is whether a runtime can
reduce end-to-end latency by selecting image scale, keypoint count, robust-fitting
budget, and, where supported by a composition-quality policy, blending mode,
while satisfying a predefined quality-of-result constraint. The planned system
combines an offline configuration search with online selection, intermediate
self-checks, and reuse of unaffected stage outputs during escalation. The
planned evaluation uses UDIS-D on an NVIDIA Tesla T4 and compares fixed
configurations, the adaptive runtime, and an offline per-pair oracle under the
same quality policy and timing boundary. No measured speedup, quality guarantee,
or completed implementation is claimed in this outline.
\end{abstract}

\section{Introduction}
\label{sec:introduction}
% Role: opening; establish the task without claiming a new stitching algorithm.
Panorama stitching combines overlapping images into a wider field of view through
alignment, warping, and composition~\cite{brown2007stitching}. This paper considers
the computational budget of a feature-based GPU pipeline rather than a replacement
for its underlying image-formation model.

% Role: challenge; state a hypothesis whose motivation needs profiling evidence.
A single fixed configuration assigns the same budgets to every image pair.
The hypothesis motivating this study is that scene-dependent quality requirements
leave opportunities to select cheaper configurations for some pairs. Whether
these opportunities outweigh selection, checking, and escalation overhead must
be established by workload characterization and held-out evaluation.

% Role: method and intended contribution; do not report planned work as results.
The proposed study asks: can a scene-aware cross-stage runtime reduce GPU
panorama-stitching latency while maintaining a predefined quality-of-result
constraint? The intended contributions are a dependency-aware budget-selection
runtime, an explicit quality-constrained configuration protocol, and a reproducible
analysis of when adaptation is beneficial. These contributions remain subject to
implementation, comparison with related runtimes, and experimental verification.

\section{Background and Problem Formulation}
\label{sec:problem}
\subsection{Pipeline and Computational Budgets}
% Role: background and notation; classical stages are not novelty claims.
The pipeline comprises preprocessing (S1), keypoint detection (S2), description
(S3), matching (S4), homography estimation (S5), and warping and blending (S6).
A configuration is $c=(s,k,r,b)$: image scale $s$ controls alignment resolution,
$k$ limits retained keypoints, $r$ limits robust-fitting hypotheses, and $b$
selects a blender. Rendering uses the original colour images. A homography
estimated in resized coordinates must be transformed back to their coordinate
systems before rendering.

\subsection{Quality Policy and Feasibility}
\label{sec:quality}
% Role: design; distinguish a relative alignment proxy from visual correctness.
The current benchmark design uses overlap SSIM and PSNR relative to a fixed
high-budget reference, together with stitch validity and coverage checks.
For pair $i$, let $\mathcal{P}_i(c)$ denote the offline feasibility predicate:
\begin{align}
\mathcal{P}_i(c) ={}& \mathrm{Valid}_i(c)\land\mathrm{CoverageOK}_i(c) \notag\\
 &{}\land[\mathrm{SSIM}_i(c)\geq\mathrm{SSIM}_{\mathrm{ref},i}-\delta_S] \notag\\
 &{}\land[\mathrm{PSNR}_i(c)\geq\mathrm{PSNR}_{\mathrm{ref},i}-\delta_P].
\label{eq:quality}
\end{align}
Development-only calibration must determine and freeze the tolerances and the
candidate-independent evaluation domain. The reference can itself fail and is
not a ground-truth panorama. These alignment proxies do not certify perceptual
composition quality or distinguish blending modes; the primary scale/keypoint/
fitting-budget experiments therefore hold $b$ fixed until an appropriate
composition policy is agreed and validated.

\subsection{Latency Objective and Offline Oracle}
% Role: objective; keep offline knowledge separate from the online decision.
For a finite configuration set $\mathcal{C}$, the offline oracle is
\begin{equation}
c_i^\star\in\mathop{\mathrm{arg\,min}}_{c\in\mathcal{C}:\mathcal{P}_i(c)}T_i(c),
\label{eq:oracle}
\end{equation}
where $T_i(c)$ uses a common end-to-end timing boundary. A pair with no feasible
configuration is recorded as infeasible. Fixed-Best is one configuration selected
on development data for all pairs; the oracle may select a different configuration
for each pair using offline quality measurements. The online controller cannot
access those measurements and must include its own overhead in reported latency.

\section{Related Work}
\label{sec:related-work}
\subsection{Stitching, Alignment, and Composition}
% Role: positioning; cite only the distinctions supported by these sources.
Classical feature-based stitching provides the pipeline context~\cite{brown2007stitching}.
UDIS++ combines global and local alignment with learned composition~\cite{nie2023udispp}.
The final review will distinguish improvements to image formation and alignment
from the runtime allocation problem studied here; it will not compare entire
stitching systems as if they differed only in their blender.

\subsection{GPU Blending and Cost--Quality Trade-offs}
% Role: prior evidence and scope; earlier GPU results are not T4 predictions.
Prior work compares panoramic blending algorithms in terms of quality, time,
and memory~\cite{zhu2018blending}. This motivates measuring blender costs and
failure modes explicitly, without transferring speedup numbers across hardware
or camera setups. Feather and multi-band modes are candidate S6 policies;
adding more pyramid levels is not assumed to improve quality monotonically.

\subsection{Adaptive Runtimes and Approximate Computing}
% Role: literature gap to complete; no unsupported first-of-its-kind claim.
This subsection will compare quality-constrained configuration selection,
feedback-based adaptation, and dependency-aware result reuse with the closest
runtime and approximate-computing literature. The contribution claim must be
narrowed after that comparison. A lack of citations in this working subsection
is an evidence gap, not evidence that no related system exists.

\section{Cross-Stage Runtime Design}
\label{sec:runtime}
\subsection{Workload Characterization and Configuration Ladder}
% Role: motivation and design; measurements must precede optimization claims.
The first experiment will measure stage costs and the per-pair feasible
configuration sets on development workloads. An offline sweep will identify
non-dominated latency--quality configurations and construct a compact ladder.
This evidence will determine whether adaptive budgets have enough headroom;
a fixed configuration or a Pareto characterization is the fallback if they do not.

\subsection{Scene-Aware Selection, Checks, and Escalation}
% Role: design; the policy is not yet an implemented or calibrated guarantee.
The runtime is planned to select an initial configuration using scene and
intermediate signals, then escalate when its calibrated checks reject the
candidate. Available S4/S5 signals include match and inlier counts, inlier ratio,
reprojection residuals, spatial coverage, and homography validity. The final
method must specify the features, decision rule, thresholds, stopping condition,
maximum retry budget, and terminal failure behaviour before test evaluation.
These proxies can make false acceptance decisions; they do not guarantee
\cref{eq:quality} on unseen pairs.

\subsection{Dependency-Aware Reuse and GPU Execution}
% Role: design and technical rationale; advantages require a reuse ablation.
Reuse is restricted to results whose inputs and conventions remain unchanged.
Increasing only the fitting budget may reuse matches and extend the hypothesis
sequence; changing scale or keypoints can invalidate downstream work. Changing
only the blender may reuse warps and validity masks if images and homography
are identical. Streams must respect stage dependencies, and materialization,
allocation, synchronization, and discarded attempts count towards policy cost.
The final implementation description will include the dependency graph and
buffer lifetimes, plus only kernel or memory optimizations supported by profiling.

\subsection{S6 Policies and Composition Boundary}
% Role: design and limitation; blend quality is a separate validation problem.
S6 uses inverse homography warping with bilinear sampling and explicit validity.
Its initial blend candidates are feather, multi-band with three reductions, and
multi-band with five reductions. A content-aware seam finder, optical flow,
and neural reconstruction are outside the planned rendering path. Blender
selection remains a separate experiment until output-sensitive quality checks
are validated. Increasing blend cost cannot be treated as a repair for an
invalid homography or arbitrary parallax.

\section{Experimental Evaluation Plan}
\label{sec:evaluation}
\subsection{Platform, Data, and Reproducibility}
% Role: protocol; no benchmark has yet produced paper-ready measurements.
The planned platform is Google Colab with an NVIDIA Tesla T4, compiling CUDA
for \texttt{sm\_75}. The evaluation must record the actual toolkit, compiler flags,
GPU state, code revision, and random seeds. UDIS-D is the primary benchmark;
development/calibration and final evaluation are separated by scene. HPatches
can support known-homography checks, and synthetic high-resolution pairs can
provide a separate resolution stress test. Neither substitutes for the primary
real-image evaluation.

\subsection{Baselines and Main Latency--Quality Comparison}
% Role: evidence plan; fixed and adaptive policies share the implementation.
Compare Fixed-Low, Fixed-High, development-selected Fixed-Best, Adaptive, and
the offline oracle using the same CUDA pipeline and timing boundary. Report
latency distributions, speedup, pair-wise quality differences, quality violations,
reference failures, and feasible-pair coverage. The main table must distinguish
execution failure from a completed stitch that violates the quality policy;
mean quality alone is insufficient.

\subsection{Stage Profiling, Overhead, and Ablations}
% Role: causal evidence plan; isolate controller and reuse effects.
Break down S1--S6 latency, controller/proxy cost, escalation attempts, data
movement, and cache materialization. Compare adaptive budgets with and without
reuse, single-stage versus cross-stage decisions, and alternative initial or
escalation policies. Hold blend mode fixed in the primary budget study.
A separate same-homography sweep will test whether S6 policies produce useful
cost--composition trade-offs before including adaptive blending in the main claim.

\subsection{Sensitivity, Scaling, and Failure Analysis}
% Role: external validity and limitations; avoid claiming multi-GPU scalability.
Evaluate sensitivity to the frozen quality tolerances, image resolution, overlap,
and available scene categories. Resolution scaling is a single-GPU workload
study, not multi-GPU scaling. Report peak workspace, dominated configurations,
small-workload overhead, parallax, motion, exposure mismatch, and cases where
no configuration is feasible. Energy measurement is optional and must include
its sampling methodology if reported.

\section{Discussion and Conclusion}
\label{sec:conclusion}
% Role: limitation; delimit the central claim before evidence exists.
This study is limited to a finite configuration set, a global-homography
rendering model, a relative alignment-quality policy, and one GPU platform.
It cannot yet claim general perceptual quality, real-time throughput, energy
savings, or superiority to state-of-the-art stitching systems. The conclusion
will answer the research question using measured end-to-end results and the
observed overhead, coverage, and quality-violation rates. If the adaptive-runtime
hypothesis is not supported, the paper will report the measured Pareto trade-off
and its limits rather than an unsupported speedup claim.

\bibliography{crossstage-stitch}
\end{document}
